\documentclass{article}
\usepackage{amsmath, amssymb, amsthm}
\usepackage{hyperref}
\usepackage{geometry}
\geometry{margin=1in}

\title{Erd\H{o}s Problem \#567}
\date{Last updated: 2025-08-31}
\author{Source: erdosproblems.com}

\begin{document}
\maketitle

\noindent\textbf{Status:} OPEN \quad \textbf{No prize}

\noindent\textbf{Tags:} graph theory, ramsey theory

\medskip

\noindent\textbf{Problem Statement:}

Let $G$ be either $Q_3$ or $K_{3,3}$ or $H_5$ (the last formed by adding two vertex-disjoint chords to $C_5$). Is it true that, if $H$ has $m$ edges and no isolated vertices, then\[R(G,H)\ll m?\]




In other words, is $G$ Ramsey size linear? A special case of [566] . In \cite{Er95} Erd\H{o}s specifically asks about the case $G=K_{3,3}$. The graph $H_5$ can also be described as $K_4^*$, obtained from $K_4$ by subdividing one edge. ($K_4$ itself is not Ramsey size linear, since $R(4,n)\gg n^{3-o(1)}$, see [166] .) Brada\'{c}, Gishboliner, and Sudakov \cite{BGS23} have shown that every subdivision of $K_4$ on at least $6$ vertices is Ramsey size linear, and also that $R(H_5,H) \ll m$ whenever $H$ is a bipartite graph with $m$ edges and no isolated vertices. This problem is #32 in Ramsey Theory in the graphs problem collection.




References

[BGS23] Brada\'C, D. and Gishboliner, L. and Sudakov, B., On Ramsey size-linear graphs and related questions . arXiv:2202.10388 (2023).

[Er95] Erd\H{o}s, Paul, Some of my favourite problems in number theory, combinatorics, and geometry . Resenhas (1995), 165-186.


\bigskip
\noindent\small{Source: \url{https://www.erdosproblems.com/567}}
\end{document}
